\section{Introduction}
\label{sec:intro}
The International Roughness Index (IRI) is the standard summary of longitudinal road roughness~\cite{sayers1986tp45,astm_e1926}.
Measuring it requires a longitudinal profile from a profilometer, so minor roads are surveyed rarely.
Production cars that already carry accelerometers could fill the gap~\cite{sattar2018review,sang2024iri,levenberg2021lira}, but their vibration depends on the vehicle, the sensor mounting and the speed, so the mapping from vibration to IRI has to be established for each car type from paired data: drives of that car over road sections whose profile was measured~\cite{sayers1986tp46}.
Collecting such paired drives is the expensive step.

Vehicle simulation offers another source of pairs.
Driving a vehicle model over profiles that road agencies already survey yields vibration with exact IRI labels, at any speed and for many vehicle parameter sets.
Estimators trained this way have been validated in simulation~\cite{jeong2020cnn,rcnn2023sim} and tested on real vehicles for the weighted longitudinal profile~\cite{nitsche2014wlp}, for profile reconstruction~\cite{ngwangwa2014exp} and for roughness classes~\cite{speedadaptive2025}; calibrated vehicle models can invert real car vibration into road profiles~\cite{skar2024inversion,agebjar2025fusion} or set an IRI--vibration relation by simulation~\cite{wu2023calib}.
What is missing is evidence, for continuous IRI from a production car on held-out roads of public data, of which parts of a simulation-trained estimator transfer, which simulator choices set its level, and how much paired real data must accompany the simulation.

We study this question on LiRA-CD~\cite{skar2023lira}, public data of one car model on four routes with profilometer references, comparing estimators trained on real cars with the same estimators trained only on simulated drives over the measured profiles of the other routes.
We vary the sensor and vehicle model of the simulator, combine it with $k$\,km of profiled road with paired drives in three ways, and compare simulated and real drives at equal profile budget, in their speed dependence on the same real segments and in their vibration spectra.

The answer has three parts.
Ranking segments within a route discriminates little: simulation-trained gradient boosting ranks them about as well as the real-trained model under every simulator variant, and an untrained, speed-normalised vibration RMS comes within \rpAbsMax{} of both on every route.
What has to be learned is the level, and it does not transfer reliably: it depends on the simulated sensor path, through the features carrying vibration above 10\,Hz, and on the simulated vehicle population.
A few kilometres of paired drives fix it: with 5\,km, simulation over all profiled roads beats training on those drives alone for every variant, whether the drives set a level offset or enter joint training; with 20\,km, joint training still does, whereas a level offset falls behind real drives alone for the variants with the largest vibration deficit at 10--15\,Hz.
Speed dependence transfers only in part, and the simulated vibration scales with speed differently from the real car's.

Our contributions are:
(i) a route-disjoint benchmark on public profilometer-referenced car data with dependence-aware inference and baselines from response-type physics (a calibrated quarter-car inversion), learned estimators and simulation-trained estimators (Section~\ref{sec:res_real});
(ii) a decomposition of what simulation transfers to a real car across nine sensor and vehicle models, with a sensor factorial, a same-segment spectral comparison and a band-limited retraining that associate the level error with the simulated vibration and test which inputs carry it (Sections~\ref{sec:res_zero} and~\ref{sec:res_mech});
(iii) the paired real data that simulation replaces, at equal profile budget and with a level offset, an affine map or joint training, for every variant, with test-route intervals and the effect of the pass density (Section~\ref{sec:res_value});
(iv) fixed-road speed elasticities of real- and simulation-trained estimators measured on real roads (Section~\ref{sec:res_speed}).
The study revises an earlier, unpublished pipeline of ours trained on a game engine; its label and protocol problems are documented in Supplement~\ref{S-sec:s_earlier}, and its corpus serves here only as a negative reference.
We will release the code, protocols, configurations and all predictions with the published version of the paper.
